\begin{proposition}[B.4.8 (Proved): a normal-line obstruction on every semistable good-reduction interface]
Let $a,b$ be holomorphic units on a disc $\Delta_q$, with
$0<|a(q)|<1$, and let $m\ne0$ be an integer.  Put
\[
 \mathcal E_a=(\C^*\times\Delta_q)/
       \bigl((z,q)\sim(a(q)z,q)\bigr).
\]
This is a smooth proper elliptic family over the full disc.
Over $\Delta_q^*$ define a line bundle $\mathcal L_m$ by
\[
 (z,v,q)\sim\bigl(a(q)z,\ q^m b(q)v,\ q\bigr).
\]
Its restriction to each elliptic fibre has degree zero.
After any finite base change $q=s^N\varepsilon(s)$, where
$N>0$ and $\varepsilon$ is a holomorphic unit, the following
extension is impossible: a proper semistable family
$p:\mathcal C\to\Delta_s$ of connected curves, with smooth
total surface, an isomorphism
$\mathcal C|_{\Delta_s^*}\simeq\mathcal E_a|_{\Delta_s^*}$,
and a line bundle on $\mathcal C$ restricting to
$\mathcal L_m$ under that isomorphism.

The assertion allows any finite collection of rational tails
in the central curve and arbitrary translations in the
punctured isomorphism.  It remains true after tensoring
$\mathcal L_m$ by the restriction of a line bundle extending
on the proposed interface.  It does not assert an obstruction
for arbitrary singular, non-flat, or higher-dimensional
interface networks.

In particular the axis $z_1=0$ in the diagonal Hopf surfaces
\[
 H_{q,a(q)}=(\C^2\setminus\{0\})/
       \langle(z_1,z_2)\mapsto(qz_1,a(q)z_2)\rangle
\]
cannot acquire such a semistable completion together with
its normal line, after any finite base change.
\end{proposition}
\begin{proof}
Choose a logarithm of $a$ on the disc and set
$\tau_a=\log a/(2\pi i)$.  Then $\Im\tau_a>0$, and
\[
 \mathcal E_a\simeq
 (\C\times\Delta_q)/(\Z+\tau_a(q)\Z).
\]
The real coordinates $z=\exp(2\pi i(x+\tau_a(q)y))$
give a smooth trivialization by the marked torus
$(x,y)\in\R^2/\Z^2$.  In particular its monodromy on
$H_1$ is the identity.  The line factors in this marking
are $1$ along the $x$ cycle and $q^m b(q)$ along the
$y$ cycle.  On a positive parameter circle, the latter
factor has winding $m$.  Equivalently, the clutching
calculation gives
\[
 \langle c_1(\mathcal L_m),[T_x]\rangle=0,
 \qquad
 \langle c_1(\mathcal L_m),[T_y]\rangle=m,
 \qquad
 \langle c_1(\mathcal L_m),[\mathcal E_{a,q}]\rangle=0.
\]
Here $T_x,T_y$ sweep the indicated cycles around the
parameter circle; the orientation of $T_y$ is chosen to
agree with the clutching convention in the displayed
line factor.  Reversing that convention changes both
sides' sign and has no effect on nonvanishing.  One can
also compute the integer on the square
$[0,1]_y\times S^1_q$: the two $y$ edges are glued by
$q^m b(q)$, whose degree is $m$.

Under $q=s^N\varepsilon(s)$ the degree becomes $mN$.
The unit $\varepsilon(s)^m b(q(s))$ has zero winding.
Thus the mixed Chern class is the nonzero covector
\[
  (0,mN)\in
 H^1(\mathcal E_{a,q};\Z)\otimes H^1(S_s^1;\Z).
\]
This is an integral calculation; no limiting logarithm
or asymptotic estimate is involved.

Suppose an extension $\mathcal C$ existed.  Its central
fibre is connected.  Flatness and properness give arithmetic
genus one, and the normalization exact sequence gives
\[
 1=\sum_i g(C_i)+b_1(\Gamma),
\]
where $C_i$ are the smooth normalization components and
$\Gamma$ is the dual graph, including a loop for a
self-node.  Consequently either there is a unique
elliptic component and $\Gamma$ is a tree, or all
components are rational and $\Gamma$ has one cycle.

The second case has nontrivial monodromy.  This can be
seen directly in the local plumbing $uv=s$.  Transport
around $s=0$ is a positive Dehn twist on its neck annulus.
An edge outside the unique graph cycle gives a separating
neck and acts trivially on $H_1$.  Each of the $k\ge1$
edges in the graph cycle gives a nonseparating neck.
These disjoint necks represent the same primitive class
up to sign on the smoothed genus-one surface.  The
transvection is unchanged by reversing that class.
Their product therefore has, in a symplectic basis, the
matrix
\[
  \begin{pmatrix}1&k\\0&1\end{pmatrix}
\]
or its inverse under the opposite monodromy convention.
It is not the identity.  The given punctured elliptic
family has identity monodromy, so this case is impossible.
This elementary plumbing argument is all that is needed
here; no algebraic stable-reduction or comparison theorem
for arbitrary analytic total spaces is being assumed.

There is therefore an elliptic core $C_0$ with rational
trees attached at finitely many points.  Choose two
oriented loops on $C_0$ avoiding those points and with
intersection number one.  The map $p$ is a submersion
near their union.  Compactness of the loops and a smooth
horizontal distribution provide, after shrinking the
disc, families of loops over the full disc.  Each swept
torus on its boundary bounds the resulting solid torus
in $\mathcal C$.  Hence the first Chern class of every
line bundle on $\mathcal C$ evaluates to zero on both
such swept tori.

The two transported loops form a basis of the nearby
elliptic fibre's $H_1$: cutting off the rational trees
adds only genus-zero pieces along separating necks,
and their intersection number remains one.  Under the
given punctured isomorphism, their swept tori therefore
project to a basis of the mixed group
$H_1(\mathcal E_{a,q};\Z)\otimes H_1(S_s^1;\Z)$.
They may differ from the marked swept tori by integral
fibre classes.  This accounts in particular for an
arbitrary loop of fibre translations in the isomorphism.
The fibre degree of $\mathcal L_m$ is zero, so these
additional fibre classes do not change either pairing.
The nonzero covector $(0,mN)$ cannot vanish on both
basis vectors.  This contradicts the bounding-solid-torus
calculation and proves the assertion.

Finally an extending tensor factor has zero pairing on
both bounding tori by the same argument.  It cannot
cancel the obstruction.  Vertical divisors and any
Cartier triple divisor on such an extended interface
are included among these harmless tensor factors.

For the stated Hopf consequence, the axis is precisely
$\mathcal E_{a,q}$.
The derivative of the deck transformation in the $z_1$
direction multiplies a normal vector by $q$; its action
on the axis coordinate is multiplication by $a(q)$.
These are exactly the two factors defining
$\mathcal L_1$.  This is also the explicit diagonal-Hopf
normal-line convention of
\citet[Section 2.2]{AbashevaDeev2024AB}.
The assertion just proved applies with $m=1$ and $b=1$.
It also applies to the first multiplier $q^m b(q)$ for
any positive $m$.  At the actual cusp the nonzero unit
$e^{2\pi i\mu(q)}$ has absolute value different from one
at $q=0$.  Indeed the nodal cubic is
$Y^2=(X-1)^2(X+2)$; its normalization
$X=h^2-2$, $Y=h(h^2-3)$ and the multiplicative coordinate
$(h-\sqrt3)/(h+\sqrt3)$ show that the marked section at
$(X,Y)=(0,\pm\sqrt2)$ has multiplier one of
$-(5\pm2\sqrt6)$.
Choose that multiplier or its inverse to obtain
such a contracting $a(q)$.
This explains why the candidate axis has good reduction.
No claim is made that an arbitrary filling is obliged
to contain this particular axis or normal factor.
\end{proof}

\begin{remark}[What a corrected cap must change]
The obstruction concerns an individual normal line,
not just the tensor product of two normal lines in
the $T^1$ condition.  Opposite windings can cancel in
that tensor product while neither line extends on its
own proposed smooth branch.  Thus cancelling the
scalar $q$ by $q^{-1}$ is insufficient for this
good-reduction interface.  A successful correction
would have to change the interface or its punctured
normal bundle outside the class specified above.
The constant-elliptic example B.4.7 has no such
winding and remains consistent with B.4.8.
\end{remark}
